\begin{tikzpicture}[font=\small,line cap=round,line join=round]\begin{scope}[x=2.25cm,y=0.25cm]\draw[otRule!65,line width=.2pt](0,0)--(0,16);\draw[otRule!65,line width=.2pt](1,0)--(1,16);\node[below,font=\tiny,fill=white,inner sep=.5pt]at(1,0){1};\draw[otRule!65,line width=.2pt](2,0)--(2,16);\node[below,font=\tiny,fill=white,inner sep=.5pt]at(2,0){2};\draw[otRule!65,line width=.2pt](3,0)--(3,16);\node[below,font=\tiny,fill=white,inner sep=.5pt]at(3,0){3};\draw[otRule!65,line width=.2pt](4,0)--(4,16);\node[below,font=\tiny,fill=white,inner sep=.5pt]at(4,0){4};\draw[otRule!65,line width=.2pt](0,0)--(4,0);\draw[otRule!65,line width=.2pt](0,2)--(4,2);\node[left,font=\tiny,fill=white,inner sep=.5pt]at(0,2){2};\draw[otRule!65,line width=.2pt](0,4)--(4,4);\node[left,font=\tiny,fill=white,inner sep=.5pt]at(0,4){4};\draw[otRule!65,line width=.2pt](0,6)--(4,6);\node[left,font=\tiny,fill=white,inner sep=.5pt]at(0,6){6};\draw[otRule!65,line width=.2pt](0,8)--(4,8);\node[left,font=\tiny,fill=white,inner sep=.5pt]at(0,8){8};\draw[otRule!65,line width=.2pt](0,10)--(4,10);\node[left,font=\tiny,fill=white,inner sep=.5pt]at(0,10){10};\draw[otRule!65,line width=.2pt](0,12)--(4,12);\node[left,font=\tiny,fill=white,inner sep=.5pt]at(0,12){12};\draw[otRule!65,line width=.2pt](0,14)--(4,14);\node[left,font=\tiny,fill=white,inner sep=.5pt]at(0,14){14};\draw[otRule!65,line width=.2pt](0,16)--(4,16);\node[left,font=\tiny,fill=white,inner sep=.5pt]at(0,16){16};\draw[otMuted,->](0,0)--(4.066666666666666,0)node[right]{h};\draw[otMuted,->](0,0)--(0,16.8)node[above]{$^\circ$C};\node[below left,font=\tiny]at(0,0){0};\begin{scope}\clip(0,0)rectangle(4,16);\draw[otNavy,line width=1pt]plot[domain=0:4,samples=90](\x,{6+2*\x});\geoSolution{}\end{scope}\end{scope}\end{tikzpicture}\par\vspace{1.5mm}{\otRegular\fontsize{8}{10}\selectfont\color{otMuted} Les axes sont gradués ; leurs échelles peuvent différer.\par}
